\documentclass[../../../../main.tex]{subfiles}
\begin{document}

\begin{flushright}
\footnotesize\textit{【自動文字起こし・要確認】}
\end{flushright}

\noindent
\textbf{[解]} 左から $k$ 番目の数の、$S_n$ での値が $S$ と一致する確率を $t_k(n)$ とする。

\bigskip
\noindent
$1^\circ$ $k=1,3,4$ の時\\
$S$ では 1 だったから、$n+1$ 回目について、$n$ 回目で場合分けして
\begin{align*}
t_k(n+1) &= (1-q) t_k(n) + p(1 - t_k(n)) \\
&= (1-p-q) t_k(n) + p \\
t_k(n+1) - \frac{p}{p+q} &= (1-p-q) \left( t_k(n) - \frac{p}{p+q} \right)
\end{align*}
$t_k(0) = 1$ から、くり返し用いて
\begin{align*}
t_k(n) &= (1-p-q)^n \left( 1 - \frac{p}{p+q} \right) + \frac{p}{p+q} \\
&= \frac{1}{p+q} \{ p + q(1-p-q)^n \} \quad \dots \text{①}
\end{align*}

\bigskip
\noindent
$2^\circ$ $k=2,5$ の時\\
$1^\circ$ と同様に、
\begin{align*}
t_k(n+1) &= (1-p) t_k(n) + q(1 - t_k(n)) \\
&= (1-p-q) t_k(n) + q \\
t_k(n+1) - \frac{q}{p+q} &= (1-p-q) \left( t_k(n) - \frac{q}{p+q} \right)
\end{align*}
$t_k(0) = 1$ から、くり返し用いて
\[
t_k(n) = \frac{1}{p+q} \{ q + p(1-p-q)^n \} \quad \dots \text{②}
\]

\bigskip
\noindent
①, ② から、
\[
C(n) = \prod_{k=1}^5 t_k(n) \xrightarrow{n \to \infty} \frac{p^3 q^2}{(p+q)^5} \quad \text{\#\#}
\]

\newpage

\end{document}
